
\chapter{Платёжный шлюз изнутри}\label{ch:gateway}

Один запрос мерчанта: \texttt{POST} к \texttt{/v1/payment\_intents}
на оплату заказа. Снаружи это обычный JSON поверх HTTPS. Внутри
же один и тот же запрос проходит через последовательность решений,
и ошибка в любом звене меняет риск, конверсию или
PCI scope. Граница между интеграцией и архитектурой проходит именно здесь.

\section{Анатомия шлюза}\label{sec:gw-anatomy}

Платёжный шлюз (payment gateway) отвечает на вопрос, что именно
должно произойти внутри системы, чтобы один запрос мерчанта
превратился в корректную, наблюдаемую авторизацию.

\definitionnote{Шлюз / процессор / оркестратор}{Три термина,
  которые часто используют как взаимозаменяемые, но на деле они
  описывают разные уровни абстракции. \textbf{Платёжный шлюз}
  (payment gateway) принимает платёжные запросы от мерчанта по
  высокоуровневому API и переводит их в протоколы карточных сетей и
  банков-эквайеров. \textbf{Платёжный процессор} (payment processor)
  непосредственно обрабатывает авторизацию, клиринг и settlement и
  часто имеет прямое подключение к карточным сетям. \textbf{Платёжный
  оркестратор} (payment orchestrator) маршрутизирует транзакции между
  несколькими шлюзами или процессорами по заданным правилам.
  Границы между ролями часто размыты: крупные провайдеры платёжных
  услуг (PSP) вроде Stripe или Adyen совмещают все три уровня в одном
  продукте.}

Независимо от названия конкретного продукта, внутри платёжного
шлюза почти всегда можно выделить пять компонентов,
через которые проходит каждая транзакция: слой приёма API-запросов,
модуль оценки риска, модуль маршрутизации, хранилище токенов
и модуль авторизации, отправляющий запрос банку-эквайеру
или напрямую карточной сети.

Порядок, в котором эти компоненты вызываются, не случаен. Он
определяет и архитектурные свойства системы, и её PCI scope.

Мерчант отправляет HTTPS-запрос на API шлюза, например,
\texttt{POST} к \texttt{/v1/payments} с JSON-телом, содержащим
сумму, валюту, данные карты или токен и метаданные заказа.
Слой приёма API-запросов валидирует сообщение, проверяет
аутентификацию (API key или OAuth token), проверяет
идемпотентность (подробнее в разделе~\ref{sec:gw-reliability})
и передаёт нормализованное внутреннее представление транзакции
следующему компоненту, модулю оценки риска.

Модуль оценки риска применяет к транзакции набор правил и моделей:
проверки частоты повторов (velocity checks), сопоставление
геолокации, отпечаток устройства, а в более зрелых системах
скоринговую модель, оценивающую вероятность мошенничества.
Результатом становится одно из трёх решений: пропустить
транзакцию дальше, отправить её на дополнительную проверку
по 3-D Secure или отклонить. Подробно устройство антифрод-систем
разобрано в главе~\ref{ch:fraud}; здесь важен сам порядок:
проверка риска работает \emph{до} авторизации, потому что
подозрительную транзакцию дешевле остановить заранее, чем потом
разбирать chargeback.

Если модуль оценки риска пропустил транзакцию, она попадает в модуль
маршрутизации,
который определяет, через какого эквайера (или какой процессор)
отправить авторизационный запрос. Логика маршрутизации остаётся одной из
самых коммерчески значимых частей
шлюза (раздел~\ref{sec:gw-routing}).

На следующем шаге, если мерчант передал полный PAN, шлюз заменяет
его токеном из хранилища токенов: PAN сохраняется в зашифрованном
хранилище, а дальше по внутреннему конвейеру и в ответе мерчанту
циркулирует только токен. Это сужает CDE (cardholder data environment)
до одного компонента и радикально сокращает PCI scope всей
системы, механизм, подробно описанный в главе~\ref{ch:pci-dss}. Если же мерчант с самого начала прислал
сетевой токен (DPAN), хранилище может не участвовать, но шлюз
всё равно проверяет ограничения домена токена (token domain
restriction) и при необходимости запрашивает у TSP актуальную
криптограмму для конкретной транзакции~\cite{emvco-payment-tokenisation-guide}
(см. также главу~\ref{ch:tokenization}).

Наконец, модуль авторизации формирует сообщение ISO~8583
(или его эквивалент в API-протоколах новых эквайеров),
отправляет запрос по защищённому каналу банку-эквайеру,
получает ответ, преобразует код ответа в формат API шлюза
и возвращает результат мерчанту. Шлюз ожидает здесь прежде
всего ответ эмитента через карточную сеть.

\begin{figure}[tbp]
  \centering
  \includefiguresvg[width=.45\textwidth]{assets/figures/ch27-payment-gateway/gateway-components}
  \caption{Компоненты платёжного шлюза.}\label{fig:gateway-components}
\end{figure}

Уже на этом минимальном маршруте видно главное свойство шлюза:
каждый следующий компонент получает более узкое и лучше
определённое представление транзакции, чем предыдущий. Именно
из этого вырастают и требования к изоляции данных, и PCI scope,
и ограничения на отказные режимы.

\section{Интеллектуальная маршрутизация}\label{sec:gw-routing}

При наличии нескольких эквайеров шлюз должен решить, куда
отправить авторизационный запрос, чтобы не проиграть
одновременно в стоимости, доле одобрений и времени ответа.

\subsection*{Три оси оптимизации}

Первая ось: \textbf{стоимость}. Разные эквайеры предлагают
разные тарифы, а interchange fee зависит от типа карты,
MCC мерчанта, региона эмитента и десятков других параметров.
Модуль маршрутизации может выбрать эквайера с минимальной
стоимостью для конкретной комбинации условий, например,
направить внутрироссийскую транзакцию по карте Мир через
эквайера с более выгодным тарифом НСПК, а трансграничную
транзакцию через другого, у которого ниже валютная
наценка.

Вторая ось: \textbf{доля одобрений}. Разные эквайеры
показывают разную долю одобрений для одних и тех же BIN-диапазонов.
Причины могут быть техническими или коммерческими:
качество подключения к сети, скорость ответа, репутация мерчанта
у конкретного эквайера, локальный или трансграничный маршрут.
Накапливая статистику по BIN, стране и сумме, модуль маршрутизации
может направить транзакцию туда, где вероятность одобрения выше.
Для крупных объёмов даже небольшая разница в качестве маршрута
быстро превращается в заметное изменение выручки.

Третья ось: \textbf{латентность}. Для мерчанта, работающего
в реальном времени, важно, чтобы ответ пришёл быстро: если
один эквайер стабильно отвечает заметно быстрее другого,
модуль маршрутизации при прочих равных предпочтёт первый,
потому что каждая лишняя секунда в сценарии оплаты
снижает конверсию.

Эти три оси неизбежно конфликтуют: самый дешёвый эквайер
может иметь худшую долю одобрений, а самый быстрый оказаться
самым дорогим. Задача модуля маршрутизации: найти рабочий
оптимум в ограничениях, заданных мерчантом: \enquote{минимизировать
  стоимость при доле одобрений не ниже X\% и p95 латентности
не более Y~мс}. Продвинутые системы решают такую задачу
как адаптивное распределение трафика; в литературе этот подход
часто описывают моделью multi-armed bandit.

\subsection*{Каскадная маршрутизация}

Каскадную маршрутизацию часто путают с интеллектуальной, хотя
она решает другой вопрос. Если первый маршрут вернул отказ,
который допускает ещё одну попытку по другому контуру, шлюз
иногда отправляет транзакцию второму эквайеру. Но термин
\enquote{soft decline} здесь нужно использовать осторожно:
у схем он может иметь узкий нормативный смысл. Например,
Mastercard требует использовать soft decline
в странах, где действует строгая аутентификация клиента
(strong customer authentication, SCA), только тогда, когда
для CNP-операции нужна, но отсутствует такая аутентификация,
а не как
общий синоним любого отказа~\cite{mc-tpr}. Логика каскада
при этом остаётся прежней: отказ одного маршрута ещё не
означает, что транзакция в принципе невозможна. Другой
эквайер может получить иной ответ от того же эмитента,
особенно если его авторизационный запрос отличается:
другой BIN эквайера, другой MCC, другое значение
Terminal ID.

\importantnote{%
Каскадная маршрутизация и обычный повтор решают разные задачи.
Каскад отправляет транзакцию \emph{другому} эквайеру
после отказа, который допускает такой манёвр; повтор же
идёт \emph{тому же} эквайеру, обычно спустя некоторое время.
Путаница между этими сценариями и приводит к дублям: если
шлюз повторил запрос тому же эквайеру, а первый ответ лишь задержался,
эмитент увидит два авторизационных запроса и может одобрить
оба. Каскад сам по себе не создаёт дубля на стороне первого
эквайера, но создаёт другой риск: задержавшийся ответ от
первого маршрута может прийти уже после того, как второй
успел получить одобрение. Поэтому грамотная реализация должна
уметь немедленно отправлять reversal первой авторизации,
если она материализовалась с опозданием.

Каскадная маршрутизация полезна только для части отказов,
которые эквайер или схема считают потенциально устранимыми.
К явно неверным реквизитам, вроде \texttt{14}: Invalid Card Number
или \texttt{54}: Expired Card, это не относится~\cite{mc-tpr}.
Кроме того, сами схемы ограничивают избыточные повторные
попытки; конкретные лимиты зависят от сценария и правил
конкретной схемы~\cite{mc-retry-rules}. Поэтому модуль
маршрутизации должен различать эти два класса отказов
по коду ответа и включать каскад только в первом случае.
}
\definitionnote{Каскад}{Повторная попытка через \emph{другого}
  эквайера или маршрут. Это не обычный retry тому же партнёру
  и не универсальный ответ на любой decline.}

Интеллектуальная маршрутизация отвечает на вопрос \enquote{куда лучше
вести первую попытку}, каскадная на вопрос \enquote{что делать после мягкого
отказа}. Смешивать эти логики опасно: у них разные цели и разные
отказные режимы.

\section{Интеграционные модели}\label{sec:gw-integration}

Способ подключения мерчанта к платёжному шлюзу определяет, где
проходит граница ответственности за карточные данные и каким будет
PCI scope мерчанта.

\definitionnote{Hosted Payment Page (HPP)}{Платёжная страница,
  размещённая на стороне шлюза. Мерчант перенаправляет туда покупателя
  и не принимает карточные данные в свой контур.}
\subsection*{Hosted payment page (перенаправление)}

В модели HPP мерчант инициирует платёжную сессию через API шлюза
и получает URL, на который перенаправляет покупателя. Покупатель
вводит данные карты на странице шлюза, шлюз обрабатывает транзакцию
и возвращает покупателя на сайт мерчанта с результатом. Параллельно
шлюз отправляет мерчанту асинхронное уведомление с финальным
статусом.

Главное преимущество HPP: минимальный PCI scope. Поскольку
карточные данные никогда не проходят через серверы мерчанта,
он обычно попадает в модель SAQ A, а не в более широкие
сценарии, где элементы платёжной страницы частично исходят
от самого мерчанта~\cite{pci-faq-1291,pci-faq-1293}. Для небольших мерчантов,
у которых нет выделенной команды безопасности, это часто
единственный реалистичный вариант.

Недостаток: потеря контроля над пользовательским сценарием:
перенаправление прерывает поток оплаты, страница шлюза может
выбиваться из визуального языка сайта, а обратный переход
после оплаты не всегда проходит гладко. Кроме того, сам сайт
мерчанта при такой схеме не исчезает из PCI scope полностью:
PCI SSC отдельно уточняет, что redirect mechanism всё равно
оставляет сайт мерчанта в области оценки~\cite{pci-faq-1332}.

\subsection*{Встроенная форма (iFrame / JavaScript SDK)}

Мерчант встраивает
в свою страницу платёжную форму, предоставляемую шлюзом
в виде iFrame или готового JavaScript-компонента.
Визуально форма
выглядит как часть сайта мерчанта, но технически данные карты
вводятся внутри iframe, контролируемого шлюзом, и отправляются
напрямую на серверы шлюза, минуя серверы мерчанта.

Эта модель сохраняет больше контроля над интерфейсом: мерчант
может стилизовать форму, не разрывать сценарий оплаты и при
этом претендовать на SAQ A или SAQ A-EP. Разница зависит
от деталей реализации и от того, проходят ли карточные данные
через JavaScript мерчанта хотя бы транзитом~\cite{pci-faq-1291,pci-faq-1293}.

\practicenote{%
В PCI DSS v4.0.1 страницы оплаты находятся под отдельным
контролем: требования 6.4.3 и 11.6.1 посвящены защите
страниц оплаты. В январе 2025 года PCI SSC обновил SAQ A,
убрав из него эти два требования и заменив их упрощённым
критерием: мерчант подтверждает, что его сайт не подвержен
атакам через скрипты. FAQ 1588 разъясняет обновлённые
критерии~\cite{pci-dss-4,pci-faq-1588}.
Даже если карточные данные вводятся в iFrame шлюза, вредоносный
скрипт на странице мерчанта может перехватить ввод через
подложный слой или keylogger ещё до того, как данные попадут
в iFrame. Поэтому при встроенной форме мерчанту
всё равно нужно доказуемо контролировать скрипты, влияющие
на payment page~\cite{pci-faq-1588}.

}
\subsection*{Прямая интеграция через API (server-to-server)}

В этой модели мерчант собирает данные карты на своей стороне
(через собственную форму) и отправляет их на API шлюза
в серверном запросе. Мерчант получает полный контроль над
интерфейсом и может реализовать произвольную предварительную
логику: проверку адреса, предавторизацию, разделение платежа
между несколькими получателями. Но платит за это
существенно более широким PCI scope: карточные данные проходят
через его серверы, и такой сценарий уже не относится к льготным
моделям SAQ A / A-EP~\cite{pci-faq-1291,pci-faq-1293,pci-saq}.

Этот путь выбирают крупные мерчанты и платёжные платформы,
которые уже имеют PCI DSS Level 1 сертификацию и для которых
прямой контроль над данными -- конкурентное преимущество
(например, возможность реализовать собственное защищённое
  хранилище токенов или оптимизировать логику повторов
на уровне PAN).

\practicenote{%
Выбор интеграционной модели -- это ещё и бизнес-решение. Разумно
начинать с iFrame или SDK: это даёт приемлемый пользовательский
опыт при минимальном PCI scope. Переходить на server-to-server
стоит только тогда, когда ограничения встроенной формы начинают
блокировать конкретный сценарий, например, если нужно
реализовать Card-on-File с прозрачной для покупателя токенизацией
или каскадную маршрутизацию на уровне PAN.
}

\section{Надёжность и наблюдаемость}\label{sec:gw-reliability}

Надёжность шлюза определяется дисциплиной обращения
с неопределённостью: ответ запаздывает, внешний вызов теряется,
мерчант нажимает \enquote{повторить}.

\subsection*{Идемпотентность}

Повторный вызов с тем же набором параметров должен давать тот же
результат, что и первый, без побочных эффектов.
Если мерчант отправил запрос на авторизацию, не получил ответа
(timeout, сетевой сбой) и повторил его с тем же idempotency key,
шлюз должен вернуть результат первой попытки, а не создавать
вторую транзакцию~\cite{stripe-idempotency}.

Для этого шлюз хранит устойчивое сопоставление между
\path{idempotency_key} и \path{transaction_id}
в хранилище с гарантированной консистентностью, обычно
в реляционной базе данных с уникальным индексом по ключу.
Срок жизни idempotency key должен перекрывать окно, в котором
ещё возможны безопасные повторы одной и той же операции:
от сетевого сбоя сразу после ответа до позднего повторного
запроса со стороны клиента.

\importantnote{%
Идемпотентность на уровне шлюза не заменяет идемпотентность
на уровне эквайера. Если шлюз отправил авторизацию эквайеру,
не получил ответа и отправил повторно, эквайер может
обработать оба запроса. Шлюз должен либо использовать
DE~11 (STAN) + DE~37 (RRN) для дедупликации
на стороне эквайера, либо реализовать сценарий
\enquote{authorize-then-check}: при timeout отправлять не повтор,
а запрос статуса у эквайера, если такой механизм предусмотрен
конкретным протоколом,
и только если статус равен \enquote{не найдено}, повторять
авторизацию.

}
\subsection*{Обработка timeout}

Timeout остаётся одной из самых неприятных ситуаций в платёжном
процессинге, потому что отсутствие ответа -- это \emph{не}
отказ. Если шлюз отправил авторизационный запрос эквайеру
и не получил ответа в течение установленного порога, возможны
как минимум три исхода: (а)~запрос вообще
не дошёл; (б)~запрос дошёл, эмитент одобрил транзакцию, но ответ
потерялся на обратном пути; (в)~запрос дошёл, эмитент отказал,
и этот ответ тоже потерялся. В момент timeout шлюз не может
различить эти три случая, и именно это делает ситуацию опасной.

Наихудший из трёх исходов: (б): деньги заблокированы на карте
держателя, эмитент считает транзакцию одобренной, а мерчант
не знает об этом и может попытаться повторить платёж, создав
двойное списание. Именно поэтому reversal остаётся штатным
способом снять неопределённость в карточном контуре: например,
Mastercard отдельно требует automatic reversal,
если после authorization request ответ не был получен~\cite{mc-tpr}.
На практике это означает, что после timeout шлюз должен сначала
попытаться закрыть неопределённое состояние и лишь затем
решать, что возвращать мерчанту.

\subsection*{Сбой шлюза между приёмом ответа и записью}

Timeout описывает ситуацию, когда ответ не пришёл. Обратная
проблема опаснее: ответ пришёл и эмитент одобрил транзакцию,
но шлюз упал до того, как персистировал результат. Состояние
транзакции расщепляется: эмитент удерживает деньги на карте
держателя, а шлюз после перезапуска видит запись со статусом
\texttt{sent\_to\_acquirer} и не знает исхода.

Если шлюз просто повторит авторизацию, эмитент получит два
hold на одну сумму -- двойное списание. Если шлюз вернёт
мерчанту ошибку, деньги останутся заблокированными без
товара. Стандартное решение -- двухфазная запись.
До отправки запроса эквайеру шлюз создаёт запись
\texttt{pending} с idempotency key и STAN. После перезапуска
фоновый процесс выбирает все \texttt{pending}-записи старше
порога (5--15~с) и выполняет status inquiry к эквайеру по STAN
и RRN. Если эквайер подтверждает одобрение, шлюз переводит
запись в \texttt{approved} и ставит webhook в очередь. Если
эквайер не знает о транзакции, шлюз записывает
\texttt{failed} и освобождает мерчанта для повторной попытки.
Если status inquiry недоступен (не все эквайеры его
реализуют), шлюз отправляет reversal, чтобы гарантированно
снять hold, и возвращает мерчанту ошибку с рекомендацией
повторить~\cite{mc-tpr}.

Без этого механизма каждый перезапуск шлюза порождает
\enquote{призрачные} hold, которые держатели карт видят как
списания, а мерчанты -- как потерянные заказы. На масштабе
тысяч транзакций в минуту даже секундный сбой создаёт десятки
таких расхождений, разрешение которых вручную через
эквайера занимает дни.

\subsection*{Дедупликация}

Идемпотентность решает лишь часть проблемы. Дедупликация
как защита от двойной обработки одной и той же транзакции
работает на нескольких уровнях. На уровне API она опирается
на idempotency key. На уровне внутреннего конвейера на
уникальность комбинации STAN + RRN + amount + timestamp,
которая отлавливает дубли, проскочившие мимо
идемпотентности, например при каскадной маршрутизации, когда
два эквайера получают запрос с разными STAN, но за одну и ту
же покупку. На уровне денежных состояний дедупликацию берёт
на себя внутренний реестр обязательств из главы~\ref{ch:internal-ledger}, а на уровне внешних файлов и банковской
выписки работает сверка из главы~\ref{ch:reconciliation}.

\subsection*{Webhooks и гарантия доставки}

Платёжный поток редко заканчивается в момент синхронного
API-ответа. Результат транзакции возвращается мерчанту
двумя путями: синхронно, в ответе на запрос, и асинхронно,
через обратный HTTP-вызов (webhook) на указанный URL.
Webhook нужен по двум причинам: синхронный ответ может
не дойти, а после первого ответа появляются
новые внешние события: завершился 3-D Secure challenge,
пришёл capture, сработал refund или поступило позднее
обновление статуса от процессора.

Публичные API уровня Stripe обычно строят доставку webhook
по принципу \enquote{как минимум один раз}: endpoint должен
быть доступен по HTTPS, система ждёт ответ 2xx, при неуспехе
повторяет доставку, а мерчант проверяет подпись события,
прежде чем доверять содержимому запроса~\cite{stripe-webhooks}.

\practicenote{%
Webhook'и у шлюза приходят как минимум один раз, поэтому обработчик
на стороне мерчанта должен быть идемпотентным: до изменения состояния
заказа нужно проверить, не обработан ли уже этот \path{event_id}.

}
\subsection*{Наблюдаемость}

При тысячах транзакций в минуту метрики должны показать деградацию
раньше, чем она затронет значительную долю трафика.
Операционная команда непрерывно отслеживает четыре группы сигналов.

\textbf{Что происходит с одобрениями.} Доля одобрений это процент
одобренных транзакций
от общего числа попыток, с разбивкой по эквайеру, BIN-диапазону,
стране эмитента, MCC и сумме. Падение доли одобрений
за короткий интервал это сигнал, требующий немедленного
расследования: причиной может быть как сбой у эквайера, так
и массовая атака с использованием скомпрометированных карт.

\textbf{Растёт ли время ответа.} p50, p95 и p99 времени обработки
запроса, измеренного от получения HTTPS-запроса мерчанта
до отправки ответа. Неожиданный рост p99 часто указывает
на проблемы с HSM, медленный ответ эквайера или перегрузку
собственного модуля маршрутизации.

\textbf{Растёт ли доля системных ошибок.} Сюда входят транзакции,
завершившиеся timeout, сетевой ошибкой или внутренней ошибкой сервера,
а не бизнес-ответом вроде approved или declined. Если эта
метрика становится заметной на фоне обычного трафика, команда
уже имеет дело с полноценным инцидентом.

\textbf{Доходят ли webhook до мерчантов.} Доля уведомлений,
доставленных с первой попытки, показывает либо проблемы
на стороне мерчантов, либо перегрузку собственной очереди
доставки.

\section{От Stripe API к ISO~8583}\label{sec:gw-api-to-iso}

Пока разработчик смотрит на транзакцию как на JSON-объект
с полями \texttt{amount}, \texttt{currency},
\path{payment_method} и \texttt{description}, банк-эквайер
и карточная сеть видят ту же самую операцию как сообщение
ISO~8583. В публичных API уровня Stripe это ещё и буквально
JSON-поверх-HTTPS с чётко описанным поведением запроса
и ответа~\cite{stripe-api-reference}. Шлюз транслирует одно
представление в другое (подробнее о самом стандарте см.
главу~\ref{ch:iso8583}).

Конвейер обработки одного \texttt{POST~/charges} внутри шлюза:

\begin{enumerate}
  \item \textbf{Приём и валидация.} HTTP-слой проверяет
    API-ключ, парсит JSON, валидирует обязательные поля
    (\texttt{amount}~>~0, \texttt{currency} по ISO~4217,
    наличие \path{payment_method}). Если передан
    \texttt{idempotency-key}, шлюз ищет кэшированный ответ
    и при совпадении возвращает его без повторной обработки.
  \item \textbf{Fraud-скоринг.} Шлюз вычисляет (или запрашивает
    у внешнего сервиса) fraud-скор. Если скор выше порога
    отказа, транзакция отклоняется до контакта с эквайером,
    экономя стоимость авторизационного запроса
    (см. главу~\ref{ch:fraud}).
  \item \textbf{Маршрутизация.} Модуль маршрутизации выбирает
    эквайера по BIN-диапазону, валюте, MCC и текущей доле
    одобрений каждого канала. Если основной канал деградировал,
    каскадное правило переключает трафик на резервный маршрут.
  \item \textbf{Извлечение PAN.} Если \path{payment_method}
    ссылается на сохранённый токен, шлюз обращается к хранилищу
    (vault) за FPAN или DPAN. Vault возвращает карточные данные
    в оперативную память процесса; на диск они не попадают
    (см. главу~\ref{ch:pci-dss}).
  \item \textbf{Сборка ISO~8583.} Шлюз формирует MTI~0100,
    заполняет bitmap и data elements согласно IFS выбранного
    эквайера. Маппинг полей описан ниже.
  \item \textbf{Отправка эквайеру.} Сериализованное сообщение
    уходит по TCP/TLS-соединению из пула. Шлюз запускает
    таймер ожидания ответа (3--10~с в зависимости от маршрута).
  \item \textbf{Парсинг ответа.} Ответ MTI~0110 парсится,
    DE~39 (Response Code) транслируется в JSON-статус. Шлюз
    персистирует результат и обновляет запись транзакции.
  \item \textbf{Ledger event.} Шлюз публикует событие
    \texttt{authorization.created} во внутреннюю шину. Модуль
    реестра обязательств фиксирует pending debit/credit
    (см. главу~\ref{ch:internal-ledger}).
  \item \textbf{Webhook.} Событие ставится в очередь доставки
    мерчанту: подписывается HMAC-ключом endpoint и отправляется
    на callback URL с ретраем при неуспехе.
  \item \textbf{Синхронный ответ.} HTTP-слой возвращает JSON
    мерчанту: \texttt{status}, \texttt{id},
    \path{decline_code} (если отказ), timestamp.
\end{enumerate}

Весь конвейер от приёма запроса до синхронного ответа укладывается
в 200--800~мс; из них 100--500~мс приходится на ожидание
эквайера и эмитента, а внутренняя обработка шлюза занимает
20--60~мс.

Маппинг JSON-полей в data elements (шаг~5).
Мерчант отправляет запрос с суммой \texttt{1500}, валютой
\texttt{RUB}, номером карты и сроком действия. Шлюз транслирует
это в сообщение MTI~0100 (Authorization Request), заполняя
десятки полей данных:

\begin{itemize}
  \item \texttt{amount} (1500) $\to$ DE~4 (Amount, Transaction):
    12-символьное числовое поле, выровненное нулями:
    \texttt{000000150000} (в минорных единицах, копейках);
  \item \texttt{currency} (\texttt{RUB}) $\to$ DE~49 (Currency Code,
    Transaction): \texttt{643} (числовой код рубля по ISO 4217);
  \item \texttt{card.number} $\to$ DE~2 (Primary Account Number):
    PAN, до 19 цифр;
  \item \path{card.exp_month}, \path{card.exp_year} $\to$
    DE~14 (Date, Expiration): \texttt{YYMM};
  \item внутренний ID терминала мерчанта $\to$ DE~41 (Card Acceptor
    Terminal Identification): 8-символьный код;
  \item MCC мерчанта $\to$ DE~18 (Merchant Type): 4-значный код.
\end{itemize}
Кроме того, шлюз генерирует поля, которых нет в API-запросе
мерчанта, но которые требуются протоколом: DE~11 (STAN),
уникальный 6-значный номер транзакции, DE~7 (Transmission Date
and Time), DE~12 (Time, Local Transaction), DE~22 (Point
of Service Entry Mode): способ ввода данных карты
и канал presentment, и DE~25 (Point of Service Condition Code).

Если мерчант использует сетевой токен вместо PAN, шлюз заполняет
DE~2 токенизированным PAN (DPAN), а криптограмму от TSP
передаёт в EMV data или в схеме-специфичном поле, зависящем
от конкретного протокола маршрута~\cite{emvco-payment-tokenisation-guide}
(см. также главу~\ref{ch:tokenization}).

Ответ эмитента проходит обратный путь: DE~39 (Response Code),
двухсимвольный код, шлюз транслирует в JSON-поле
\texttt{status} (\texttt{succeeded}, \texttt{declined},
\path{requires_action}) и дополняет человекочитаемым
описанием. При этом шлюз неизбежно \emph{теряет} часть
информации: ISO~8583 различает десятки причин отказа (\texttt{05}:
  Do Not Honor, \texttt{51}: Insufficient Funds, \texttt{61}: Exceeds
Withdrawal Limit), тогда как API шлюза обычно группирует их
в несколько категорий. Разработчику, которому важно отличить
\enquote{недостаточно средств} от \enquote{карта заблокирована},
нужно смотреть на поле \path{decline_code}
(или его аналог) в дополнение к \texttt{status},
подробнее об интерпретации кодов ответа см.~\cite{mc-tpr}
в главе~\ref{ch:decline-management}.

\practicenote{%
Многие
шлюзы предоставляют подробный протокольный лог транзакции
(raw response), в котором видны значения data
elements: DE 39,
DE 44 (Additional Response Data), DE 55 (EMV data). Если шлюз
этого не показывает и эквайер возвращает непонятный отказ,
запросите у эквайера trace log по STAN и RRN: это единственный
способ понять, что именно произошло на уровне протокола.
}

Трансляция охватывает авторизацию и другие операции. Refund, void
и capture не сводятся к одному универсальному набору сообщений
для всех шлюзов и эквайеров: конкретное соответствие зависит
от стадии операции и протокола конкретного маршрута. На уровне API это выглядит как несколько продуктовых
действий, а на уровне карточного контура как последовательность
authorisation, reversal, financial messages и clearing updates,
которые потом оказываются во внутреннем реестре обязательств
и в сверке. Поэтому сам шлюз лучше понимать как пограничный
слой. Окончательная правда о деньгах живёт во внутренних реестрах и
расчётных системах.
